\subsection{Primary multi-seed protocol}

The primary evidence comprises two independent full-protocol runs that differ only by random seed.
Each run uses Meta Llama~3.1~8B with temperature $0.5$ and $\mathrm{top}\_p=0.95$ for ordinary decisions, lower temperatures for pairwise scoring and voting, and a 4096-token context window.
Each society contains 26 agents (12 developed, 14 developing) interacting for 30 rounds.
Institutional assignment is forced every round.
Developed agents enter SI.
Developing agents enter SFI.
Public-good multiplier $m$ equals 1.6.
Climatic shocks occur in rounds~5 and~10 with severities 0.10 and 0.20.
Democracy sessions occur every five rounds and therefore coincide with both shocks.
Vulnerability is 0.30 for developed agents and 1.50 for developing agents.
Stage-2 budgets in SI equal $\max(20,\lfloor 0.05\times\mathrm{wealth}\rfloor)$.
Default punishment effect-to-cost is 3; reward effect-to-cost is 1.
Gossip bulletins are eligible when pairwise scores are at most 7, with a capacity of five items.

Both runs are analyzed with equal weight.
We do not designate either seed as primary and the other as supplemental.
Where signs agree, we treat the association as robust under this protocol.
Where signs diverge, we report heterogeneity rather than pooling into a single universal effect.
Agent identifiers are not treated as fixed personas across seeds, because group labels reshuffle for a substantial subset of identifiers.

\subsection{Outcomes and event definitions}

The primary outcome is $\mathrm{prop}_{i,t}$.
Immediate event-study changes use $\Delta\mathrm{prop}_{i,t}=\mathrm{prop}_{i,t+1}-\mathrm{prop}_{i,t}$.
Bad reputation uses aggregate peer-trust below 4.
Reputation drops use declines of at least 1.0.
Gossip targets follow reconstructed eligibility consistent with the bulletin rule.
Diff-in-diff contrasts, when reported for a seed, compare affected agents with same-round co-institution controls without the focal event.

Secondary outcomes include zero-contribution shares, round-mean trajectories by institution, Mann--Kendall trend tests on round means, developed--developing wealth gaps, loss-and-damage coverage ratios, adopted democratic rules, and coded motifs in post-event rationales.
SI versus SFI mean comparisons are descriptive.
Because institution membership is collinear with group label under forced routing, those comparisons are not causal institution effects.

\subsection{Qualitative coding}

For the first seed, contribution rationales in rounds immediately after bad-reputation or gossip events are coded with a multi-label codebook.
Opportunistic covers free-riding, self-interest, payoff maximization, or marginal-return language.
Reputation management covers explicit standing, credibility, or repair language.
Conformity covers peer-average imitation.
A second coder independently labels a stratified subsample for reliability.
Cross-seed language dialects are summarized with keyness contrasts between SI and SFI contribution rationales.

\subsection{Boundary baselines}

Shock-free free-choice baselines from related abstract public-goods runs motivate the climate design.
In those baselines, reputation can pull agents toward SI when exit is available, and greedy zero contribution dominates without existential risk.
We cite those patterns only as boundary conditions.
They are not re-estimated under the forced climate protocol and are not pooled with the two climate seeds.
